\section{Analysis of \sys{}}
\label{sec:analysis}

This section shows the properties of \sys{}. All results use the model of \cref{sec:model}:
not more than $f$ officers are hostile, $n \ge 3f + 1$, and
$q = \lceil (n + f + 1)/2 \rceil$. The attacker controls all links in all results. Only for
\cref{thm:tolerance}, it is necessary that messages come through.

\subsection{Two Lemmas}

\begin{lem}[One grant for each epoch]
\label{lem:grant}
A run has not more than one valid grant for each epoch.
\end{lem}

\begin{proof}
Let a run have two different valid grants of an epoch $j > 1$. Each grant has $q$ signers,
thus one loyal officer $u$ signed the two grants. By rule O2, $u$ signs a grant of epoch
$j$ only if $j > E_u$, and then $u$ sets $E_u$ to $j$. The number $E_u$ does not decrease.
Thus $u$ signs not more than one grant of epoch $j$. This is a contradiction.

For epoch 1, the parties have one valid grant at the start. A second valid grant of epoch 1
has a loyal signer, because $q > f$. But a loyal officer signs no grant of epoch 1, because
$E_u \ge 1$ at all times.
\end{proof}

\begin{lem}[Separation]
\label{lem:separation}
Let $\renew{k,e}$ and $\grant{j,h,b,s}$ be valid, with $j > k$. Then $s \ge e + F + g$.
\end{lem}

\begin{proof}
The renewal and the grant have $q$ signers each, thus one loyal officer $u$ signed the two
certificates. First, let $u$ sign the renewal before the grant. Then $M_u \ge e$ when $u$
signs the grant, and rule O2 gives $s \ge M_u + F + g \ge e + F + g$. Second, let $u$ sign
the grant before the renewal, thus $E_u \ge j > k$ from that time. By rule O1, $u$ signs
the renewal only if $k = E_u$. This is a contradiction.
\end{proof}

\subsection{Unity of Command}

\begin{thm}[Unity of command]
\label{thm:unity}
Let $F = 0$ and $g > 2\eps$. Then \sys{} has unity of command with separation
$g - 2\eps$.
\end{thm}

\begin{proof}
Let a device $d$ accept an order of an officer $x$ at real time $t$. Let a device $d'$
accept an order of a different officer $y$ at real time $t'$. Rule D2 does not apply,
because $F = 0$. By rule D1, the orders are for epochs $k$ and $j$, and $x$ and $y$ are the
commanders in valid grants of these epochs. By \cref{lem:grant}, $k \ne j$, because
$x \ne y$. Let $k < j$.

The device $d$ has a valid $\renew{k,e}$ with $C_d(t) \le e$, thus $t \le e + \eps$. The
device $d'$ has a valid $\grant{j,y,b,s}$ with $s \le C_{d'}(t')$, thus $t' \ge s - \eps$.
By \cref{lem:separation}, $s \ge e + g$. Thus $t' - t \ge g - 2\eps$. If $k > j$, the same
argument gives $t - t' \ge g - 2\eps$.
\end{proof}

The theorem applies with a hostile commander, with hostile officers, and with each sequence
of messages. It tells that all orders of one epoch come before all orders of a higher
epoch.

\subsection{Tenure}

\begin{thm}[Tenure of a commander]
\label{thm:tenure}
Let $k$ be an epoch and $t_0$ a real time. Let $W$ be a set of $n - q + 1$ or more loyal
officers. From $t_0$, let no officer in $W$ sign a renewal of epoch $k$. Then no device
accepts an order of the commander for epoch $k$ at a real time $t \ge t_0 + H + 2\eps$.
Also, no device accepts an order of the deputy for epoch $k$ at a real time
$t \ge t_0 + H + F + 2\eps$.
\end{thm}

\begin{proof}
Let $\renew{k,e}$ be a valid renewal. It has $q$ signers, and not more than $q - 1$
officers are not in $W$. Thus one signer $u$ is in $W$. The officer $u$ is loyal, and it
signed at a real time $t_u < t_0$. By rule O1, $e \le C_u(t_u) + H \le t_u + \eps + H$.
Thus $e < t_0 + H + \eps$ for each valid renewal of epoch $k$.

Let a device $d$ accept an order by rule D1 at real time $t$. Then $C_d(t) \le e$ for a
valid renewal, thus $t \le e + \eps < t_0 + H + 2\eps$. Let $d$ accept an order by rule D2
at real time $t$. Then $C_d(t) \le a + F$, and the anchor $a$ is the end time of a valid
renewal. Thus $t < t_0 + H + F + 2\eps$.
\end{proof}

The proof uses no property of the network. The attacker cannot extend the tenure when it
stops, delays or replays messages. It can only make the tenure shorter.

\begin{cor}[Two thresholds]
\label{cor:thresholds}
Let $r_h = n - q + 1$ and $r_l = n - q + 1 - f$.
\begin{enumerate}
  \item If $r_h$ loyal officers stop their consent at $t_0$, the tenure of the commander
  stops before $t_0 + H + 2\eps$. The hostile officers cannot prevent this.
  \item The hostile officers cannot make a valid certificate without $q - f$ loyal signers.
  \item If the number of loyal officers who stop consent is less than $r_l$, then $q$ or
  more loyal officers continue to sign. Thus the renewals continue when the hostile
  officers do not sign.
\end{enumerate}
\end{cor}

\begin{proof}
The first part is \cref{thm:tenure}. The second part is correct because $q > f$. For the
third part, the number of loyal officers is $n - f$ or more. If less than $r_l$ of them
stop, more than $n - f - r_l = q - 1$ continue.
\end{proof}

The two thresholds are different for the two types of commander. To remove a hostile
commander, $r_h$ loyal officers are sufficient. To stop the renewals of a loyal commander,
$r_l$ loyal officers must stop consent, together with $f$ hostile officers. For $n = 3f + 1$
the values are $q = 2f + 1$, $r_h = f + 1$ and $r_l = 1$. Thus one loyal officer and $f$
hostile officers can stop the renewals of a loyal commander. A larger $n$ makes $r_l$
larger: $n = 7$ and $f = 1$ give $q = 5$, $r_h = 3$ and $r_l = 2$.

The numbers of officers for removal and for a new commander are different. For removal,
$r_h$ loyal officers are sufficient. For a new grant, a quorum is necessary. If $r_h$ officers stop
consent and no quorum signs a grant, the system has no commander. Then only local control
stays (rule D3).

\subsection{Outage Tolerance}

\begin{thm}[Outage tolerance]
\label{thm:tolerance}
Let all officers be loyal and give consent to the commander $c$ of an epoch $k$. Let each
officer $u$ have $E_u = k$ and the valid grant of epoch $k$, with the start time $s$. Let
$c$ start a request at intervals of $\Delta$ or less in real time. Let $c$ have the valid
renewal not more than $\delta$ after the start of each request.

Let an outage start at a real time $t_0 \ge s + \Delta + \delta + \eps$. From $t_0$, $c$
and a device $d$ receive only messages that $c$ or $d$ sent, and the switch of $d$ is in
the position ``REMOTE''. Then $d$ accepts each order of $c$ that it receives at a real time
$t$ with $t_0 \le t \le t_0 + H - \Delta - \delta - 4\eps$.
\end{thm}

\begin{proof}
One request started at a real time $t_r$ with $t_0 - \Delta - \delta < t_r \le t_0 -
\delta$. Its end time is $e = C_c(t_r) + H - 2\eps$, thus $e \ge t_r + H - 3\eps$ and
$e \le t_r + H - \eps$. An officer $u$ who receives the request at a real time
$t_u \ge t_r$ has $C_u(t_u) + H \ge t_r - \eps + H \ge e$. Thus each officer can sign by
rule O1, and $c$ has the valid renewal at $t_0$.

Each order of $c$ has the grant and the renewal, and rule D3 does not stop the order. The
device accepts by rule D1 at $t$ if $s \le C_d(t) \le e$. The first condition is
correct because $C_d(t) \ge t_0 - \eps \ge s$. The second condition is correct if
$t + \eps \le e$. This is correct for each $t \le t_r + H - 4\eps$, and
$t_r > t_0 - \Delta - \delta$.
\end{proof}

Thus the outage tolerance of \sys{} is $H - \Delta - \delta - 4\eps$ or more. By
\cref{thm:tenure}, the hostile tenure is not more than $H + 2\eps$. The difference between the
two quantities is not more than $\Delta + \delta + 6\eps$. By \cref{thm:limit}, the
difference cannot be less than zero for a protocol. With the values of \cref{tab:params},
the outage tolerance is 567.2\,s or more. The hostile tenure is not more than 600.4\,s, and the
difference is not more than 33.2\,s.

\subsection{Fallback to the Deputy}

\begin{prp}[Fallback]
\label{prp:fallback}
Let $F > 0$ and $g > 2\eps$. Then \sys{} has these properties.
\begin{enumerate}
  \item At one clock time, one device accepts orders of not more than one officer.
  \item For each epoch, one device accepts orders of the deputy only at clock times in one
  interval of length $F$.
  \item Let a device accept an order of epoch $k$ at real time $t$. Let a device accept an
  order of an epoch $j > k$ at real time $t'$. Then $t' - t \ge g - 2\eps$.
\end{enumerate}
\end{prp}

\begin{proof}
For the first part, let the two orders be for the same epoch. By \cref{lem:grant}, one
order is of the commander and one order is of the deputy. By rule D1, the device has a
valid renewal with an end time $e \ge \tau$. By rule D2, the device has no such renewal.
This is a contradiction.

Then let the orders be for the epochs $k < j$. For the order of epoch $j$, the start time
is $s \le \tau$. For the order of epoch $k$, $\tau \le e + F$ for the end time $e$ of a
valid renewal. By \cref{lem:separation}, $s \ge e + F + g > \tau$. This is a
contradiction.

The second part is correct because the anchor $a$ of an epoch does not change after the
first order of the deputy. Rule D2 accepts an order only if $a < \tau \le a + F$. For the third part, the
order of epoch $k$ gives $t \le e + F + \eps$ for the end time $e$ of a valid renewal. The
order of epoch $j$ gives $t' \ge s - \eps$. By \cref{lem:separation}, $s \ge e + F + g$.
\end{proof}

With $F > 0$, \sys{} does not have unity of command. A device with no new renewal can
accept orders of the deputy while a different device accepts orders of the commander. The
proposition puts a limit on this condition: one interval of length $F$ for each device and
each epoch. The deputy exposure at a device is thus not more than $F + 2\eps$ in real time
for each epoch. By \cref{thm:deputy}, a protocol with the deputy period $F$ cannot have a
smaller exposure.
